\part{分布式训练}

\chapter{并行策略}
\label{sec:distributed}

\begin{motivationbox}{单卡容量和单卡吞吐都无法继续扩展}
模型变大后会依次遇到两类限制：模型虽然能放入单卡，但训练速度太慢；或者参数、优化器状态与激活根本放不下。分布式训练不是单一算法，而是根据瓶颈选择切分轴：切批次提高吞吐，切矩阵或层容纳模型，切序列降低长上下文激活，切专家承载 MoE 参数。
\end{motivationbox}

大模型训练的核心问题是：参数、梯度、优化器状态、激活和批次都可能放不下单卡。不同并行策略切分的对象不同。

\begin{table}[htbp]
\centering
\small
\begin{tabularx}{0.98\textwidth}{lYYYY}
\toprule
方法 & 切分对象 & 通信 & 显存收益 & 典型用途\\
\midrule
数据并行（DDP） & 批次 & 梯度全归约 & 参数仍复制 & 中小模型多卡训练\\
张量并行 & 矩阵维度 & 全归约/全收集 & 切参数与激活 & 单层太大\\
流水线并行 & 层 & 发送/接收激活 & 切模型层 & 层数很多\\
序列并行 & 序列维度 & 全收集/归约散射 & 降低激活 & 长上下文\\
专家并行 & 专家 & 全交换 & 分散专家参数 & MoE\\
ZeRO-1 & 优化器状态 & 归约与收集 & 节省优化器状态 & Adam 状态大\\
ZeRO-2 & 再切梯度 & 归约散射 & 进一步节省 & 大规模稠密模型\\
ZeRO-3/FSDP & 再切参数 & 全收集参数 & 最大节省 & 超大模型\\
\bottomrule
\end{tabularx}
\caption{常见分布式训练方法。}
\end{table}

\section{为什么需要不同的并行维度}

\textbf{数据并行（Data Parallel）与分布式数据并行（DDP）}适用于完整模型能放入单卡，但希望用更多样本并行提高吞吐的情况。每张卡计算不同数据，代价是反向后必须同步梯度。它不切模型参数，因此不能解决“单卡放不下模型”。

\textbf{张量并行（Tensor Parallel，TP）}针对单个 Transformer 层或大矩阵本身放不下单卡的问题，沿隐藏维、注意力头或前馈网络中间维切分矩阵乘法。每层前向和反向都可能需要全归约（AllReduce）或全收集（AllGather），通信频繁，所以通常优先放在节点内高速互联的 GPU 之间。

\textbf{流水线并行（Pipeline Parallel，PP）}针对模型层数多、整体参数放不下的问题，把连续层分配到不同阶段（stage）。它只在阶段边界传激活，但后一个阶段必须等待前一个阶段的输出。微批处理用流水线重叠不同样本来减少空闲，仍会留下流水线气泡，也会增加调度复杂度。

\textbf{序列并行（Sequence Parallel）}针对长上下文下激活随 $T$ 增长的问题，沿序列维分散层归一化、随机失活或部分注意力/前馈网络激活。它节省的是激活而非模型总参数，并需要在进入依赖完整隐藏维或序列信息的算子前做集合通信。

\textbf{专家并行（Expert Parallel）}针对 MoE 的专家总参数过大，把不同专家放到不同设备。路由器决定词元去向后执行全交换分发，再将专家输出合并回原位置。其效率高度依赖负载均衡和网络带宽。

实际大模型常组合数据并行、张量并行、流水线并行与 ZeRO/FSDP，形成多维并行。选择顺序应从瓶颈出发：先问模型是否放得下，再问激活是否放得下，最后评估哪种通信能被当前硬件拓扑承受。

\begin{table}[htbp]
\centering
\small
\begin{tabularx}{0.98\textwidth}{lYYYY}
\toprule
方法 & 切了什么 & 没切什么 & 主要通信 & 主要节省\\
\midrule
DDP & 批次 & 参数、梯度状态副本 & 梯度 AllReduce & 提高吞吐，不省模型状态\\
TP & 单层矩阵 & 通常不切层数 & 层内 AllReduce/AllGather & 参数与部分激活\\
PP & 连续层 & 单层内部矩阵 & 阶段间激活发送/接收 & 每阶段模型层\\
SP & 序列维激活 & 通常不切全部参数 & AllGather /\newline ReduceScatter & 长序列激活\\
EP & 专家集合 & 非专家主干 & 词元 All-to-All & 专家参数\\
ZeRO/FSDP & 模型状态 & 视阶段而定的激活 & 参数收集与梯度归约 & 参数/梯度/优化器状态\\
\bottomrule
\end{tabularx}
\caption{从“切了什么、没切什么”理解分布式策略。}
\end{table}

\begin{interviewbox}{为什么数据并行与张量并行经常组合？}
\textbf{20 秒回答：}张量并行解决单个模型副本或单层放不下，把一层矩阵切到设备组内；数据并行再复制这个 TP 设备组，处理不同批次提高吞吐。它们切的是不同轴，因此可以组成二维并行。

\textbf{深入解释：}TP 每层通信频繁，通常放在节点内高速互联；DP 的梯度同步粒度更粗，可以跨节点。只用 DP 时每卡仍需完整模型，只用 TP 时虽然模型放得下，但全局 batch 吞吐未必足够。组合方式应匹配网络拓扑，而不是只追求更多并行维度。

\textbf{面试官继续追问：}
\begin{enumerate}
  \item TP 为什么几乎每层都有通信？
  \item PP 的 bubble 从哪里产生？
  \item micro-batch 为什么能减少 bubble？
  \item EP 为什么与 MoE 强绑定？
\end{enumerate}
\end{interviewbox}

\section{分布式数据并行与全归约}

\begin{motivationbox}{多卡计算必须保持同一份模型参数}
把不同小批次（mini-batch）分给多张卡可以近似线性增加样本吞吐，但各卡会得到不同局部梯度。若直接更新，模型副本从第一步起就会分叉。DDP 通过全归约求全局平均梯度，使多卡结果在数学上等价于一个更大的批次。
\end{motivationbox}

DDP 中每张卡持有完整模型，处理不同小批次分片。反向传播后，第 $i$ 张卡得到局部梯度 $g_i$。设设备总数为 $N$，为了等价于大批次训练，需要同步平均：
\[
g=\frac{1}{N}\sum_{i=1}^{N}g_i.
\]
其中 $g$ 是所有设备的平均梯度。全归约完成“所有设备求和并把结果发回所有设备”；如果不同步梯度，各卡参数会向不同方向更新。

这里不能用普通 Gather 代替：Gather 只是把各卡张量拼到某个目标设备，既没有求和，也不会把结果自动发回所有设备；AllReduce 同时执行归约和广播语义，使每张卡都获得同一全局梯度。实现上可以分解为 ReduceScatter 与 AllGather，但语义仍是“全局归约后每卡得到结果”。

\begin{examplebox}{两张卡为什么必须平均梯度}
同一参数在卡 1 的局部梯度为 $g_1=(2,0)$，卡 2 为 $g_2=(0,4)$。AllReduce 求和后两张卡都得到 $(2,4)$，再除以设备数得到全局平均
\[
g=(1,2).
\]
若学习率为 $0.1$，两张卡都会执行同样的参数更新 $-0.1(1,2)$。若不同步，卡 1 更新 $(-0.2,0)$、卡 2 更新 $(0,-0.4)$，下一步两份模型已不再相同，也就不能继续表示同一个大批次训练。
\end{examplebox}

对大小为 $S$ 字节的梯度，环形全归约（ring AllReduce）可分解为归约散射（ReduceScatter）与全收集。每阶段通信 $(N-1)S/N$ 字节，因此每张卡总通信量近似为
\[
V_{\rm ring}=2\frac{N-1}{N}S.
\]
其中 $V_{\rm ring}$ 是每张卡发送和接收的数据量，$N$ 是设备数。卡数增大时它趋近 $2S$，所以 DDP 扩展效率取决于反向计算能否覆盖这部分通信，而不只是 GPU 数量。

\section{ZeRO 与 FSDP 的内存推理}
\label{sec:training-memory}

\begin{motivationbox}{DDP 的每张卡为什么仍然很快耗尽显存？}
DDP 只切数据，每张卡仍复制完整参数、梯度和 Adam 状态。模型参数规模增长时，增加 GPU 数量并不会降低单卡模型状态显存。ZeRO/FSDP 把这些原本重复的状态分片存放，在计算需要时临时聚合，用更多通信换取近似 $1/N$ 的状态显存。
\end{motivationbox}

Adam 训练常存参数、梯度、一阶矩、二阶矩。若每项都按 FP32 估算，优化器状态可比参数本身还大。ZeRO/FSDP 的思路是把这些状态分片（shard）到不同设备，需要时再收集（gather）。代价是通信增加，收益是单卡显存显著下降。

\begin{mathbox}
以混合精度 Adam 为例，每个参数常近似需要：BF16 参数 $2$ 字节、BF16 梯度 $2$ 字节、FP32 主参数副本 $4$ 字节、两个 FP32 动量状态 $8$ 字节，合计约 $16$ 字节。若模型有 $P$ 个参数，设备数为 $N$，忽略激活与临时缓冲区：
\[
M_{\rm DDP}\approx16P,
\]
\[
M_{\rm ZeRO1}\approx4P+\frac{12P}{N},\qquad
M_{\rm ZeRO2}\approx2P+\frac{14P}{N},\qquad
M_{\rm ZeRO3}\approx\frac{16P}{N}.
\]
这些公式中的 $M$ 表示单卡持久模型状态的字节数。它们是用于面试估算的理想值；真实峰值还包含激活、通信桶（bucket）、参数全收集和显存碎片。
\end{mathbox}

\begin{examplebox}{10 亿参数在 4 张卡上的理想状态显存}
令 $P=10^9$，按每参数约 $16$ 字节估算。DDP 在每张卡都保存完整状态，因此约需 $16$ GB。理想 ZeRO-3/FSDP 将参数、梯度和优化器状态都均分到 $N=4$ 张卡：
\[
M_{\rm ZeRO3}\approx\frac{16P}{4}=4\text{ GB/卡}.
\]
节省的 $12$ GB 并没有消失，而是分散到了其他设备；某层计算前仍要临时全收集它的参数。所以 $4$ GB 是持久模型状态估算，不等于训练峰值显存。
\end{examplebox}

ZeRO 的三个阶段可以理解为逐步消除 DDP 的重复副本：阶段 1 切优化器状态，阶段 2 再切梯度，阶段 3 连参数也切。完全分片数据并行（Fully Sharded Data Parallel，FSDP）与 ZeRO-3 的核心目标相近，都是在层计算前全收集所需参数、计算后再释放或重新分片。分片越彻底，显存越省，但参数通信和调度越敏感。

\subsection{训练显存的完整核算框架}

面试中先写出
\[
M_{\rm train}=
M_{\rm param}+M_{\rm grad}+M_{\rm opt}
+M_{\rm act}+M_{\rm temp}+M_{\rm comm}.
\]
六项分别是参数、梯度、优化器状态、前向激活、算子临时缓冲和通信缓冲。模型权重只是第一项；序列越长、微批次越大或未使用高效 attention 时，激活与临时缓冲也可能成为峰值主导项。

推理通常不需要参数梯度、Adam 一二阶动量和反向激活，因此比训练少很多状态；但长上下文服务会新增键值缓存。梯度累积只降低单次微批次激活，不会减少参数、梯度或优化器状态。激活检查点用重复前向计算换激活显存；ZeRO/FSDP 主要分片模型状态，二者解决的不是同一项，常组合使用。

\begin{exercisebox}{为什么 7B BF16 模型不只用 14 GB 就能训练？}
7B 参数的 BF16 权重约为
\[
7\times10^9\times2\text{ bytes}=14\text{ GB}.
\]
若采用前述每参数约 16 字节的混合精度 Adam 简化模型，单卡持久模型状态约为
\[
7\times10^9\times16\text{ bytes}=112\text{ GB}.
\]
这还没有包含激活、临时张量、通信桶和碎片。若理想 ZeRO-3 分到 8 张卡，持久状态约为 $112/8=14$ GB/卡，但层参数 AllGather 与激活会使真实峰值高于 14 GB。
\end{exercisebox}

\begin{interviewbox}{为什么训练显存远大于推理显存？}
\textbf{20 秒回答：}训练除权重外，还要保存梯度、Adam 状态和用于反向的中间激活，并存在临时与通信缓冲；推理通常只保留权重、当前激活和 KV Cache。因此“7B BF16 权重 14 GB”只是推理权重下界，不是训练显存。

\textbf{深入解释：}混合精度 Adam 常保留 BF16 权重/梯度、FP32 主权重和两个 FP32 动量，约 16 字节/参数。ZeRO 分片模型状态，checkpointing 压缩激活，gradient accumulation 只改变单步激活与逻辑 batch；三者不能互相替代。

\textbf{面试官继续追问：}
\begin{enumerate}
  \item gradient accumulation 是否降低 parameter memory？
  \item activation checkpointing 用什么换显存？
  \item ZeRO-2 与 ZeRO-3 多切了什么？
  \item inference 中什么状态可能随上下文增长？
\end{enumerate}
\end{interviewbox}

\begin{interviewbox}{FSDP forward 前为什么需要 AllGather？}
\textbf{20 秒回答：}持久状态中每张卡只保存当前层参数的一个分片，但完整矩阵乘法需要该计算组内的完整参数。因此进入层前临时 AllGather 参数，计算后再释放或重新分片，以通信换单卡持久显存。

\textbf{深入解释：}ZeRO-1 只分片优化器状态，ZeRO-2 再分片梯度，ZeRO-3/FSDP 连参数也分片。分片越彻底，持久显存越低，但参数预取、通信重叠和峰值临时内存越关键；理想 $1/N$ 不能直接当作真实峰值。

\textbf{面试官继续追问：}
\begin{enumerate}
  \item 参数 AllGather 何时可以与计算重叠？
  \item 为什么 bucket 太大或太小都可能不好？
  \item FSDP 与 TP 都切参数，区别是什么？
\end{enumerate}
\end{interviewbox}

\section{激活检查点（Activation Checkpointing）}

\begin{motivationbox}{模型状态能放下，前向激活仍可能放不下}
反向传播需要中间激活来计算梯度，深层网络和长序列会保存大量 $[B,T,C]$ 张量。这里 $B$ 是批次大小，$T$ 是序列长度，$C$ 是隐藏维度。ZeRO/FSDP 主要切模型状态，并不能消除这些激活。激活检查点只保存少数边界激活，反向时重新执行部分前向，以额外计算换显存。
\end{motivationbox}

若普通训练保存 $L$ 层全部激活，激活显存大致随层数 $L$ 线性增长。分段检查点可把若干连续层视为一段，只保存段输入；反向到该段时重算内部激活。分段数量越少，显存越省但重算越多。它不减少参数、梯度或优化器状态，也不会直接减少通信，因此通常与 FSDP、混合精度和序列并行组合使用。

\begin{examplebox}{用重算换掉哪些激活}
假设网络有 12 层。普通训练保存每层输出 $h_1,\ldots,h_{12}$，其中 $h_\ell$ 表示第 $\ell$ 层的输出；若每份激活占 1 GB，仅这些边界张量就约 12 GB。现在每 4 层设一个检查点，只长期保存输入 $h_0,h_4,h_8$ 和最终输出。反向计算第 9--12 层时，先从 $h_8$ 重新前向得到内部激活，再立即用于反向并释放。

实际峰值还包含当前分段的临时激活，因此不会精确降到 3 GB；但例子说明了交换关系：保存的长期激活减少，部分前向计算会在反向阶段重复执行。
\end{examplebox}

分布式训练解决的是“如何把固定训练任务放到多设备”。部署后的知识时效性与外部操作能力仍不是扩大并行度能解决的问题，这分别引出检索增强生成和智能体。

\section*{面试高频问题}
\addcontentsline{toc}{section}{面试高频问题：分布式}

\textbf{Q: 张量并行和流水线并行的区别？}

20 秒回答：TP 切单层矩阵，解决一层放不下；PP 切模型层，解决整体层数太多。TP 通信更频繁，PP 要处理流水线气泡。

\textbf{Q: ZeRO-3 为什么比 ZeRO-2 更省显存？}

ZeRO-2 切优化器状态和梯度；ZeRO-3 还切参数本身，因此参数也不再每卡完整复制。
